% !TEX program = xelatex
\documentclass[11pt, UTF8, a4paper]{ctexart} % 如果内容很多，可以改为 ctexrep

% 宏包引入
\usepackage[margin=2.5cm]{geometry} % 页面边距
\usepackage{graphicx} % 插入图片
\graphicspath{{images/}}
\usepackage{amsmath, amssymb} % 数学公式（运动学、PID计算必备）
\usepackage{hyperref} % 超链接与书签
\usepackage[table]{xcolor} % 颜色支持
\usepackage{listings} % 代码块排版
\usepackage{tabularx}
\usepackage[most]{tcolorbox} % 提示框（用于写“踩坑提示”）
\usepackage{array}
\usepackage{booktabs}
\usepackage{twemojis}
\usepackage{listings}
\usepackage{menukeys}
\usepackage{float} % 强制图片位置

% 样式设置
% 超链接样式
\hypersetup{
    colorlinks=true,          % 开启链接颜色（去掉难看的方框）
    urlcolor=black,   % 网页链接的颜色（优雅的深蓝色）
    linkcolor=black,  % 内部链接（目录、章节）的颜色（黑色）
}

% 定义代码高亮使用的颜色
\definecolor{codegreen}{rgb}{0,0.6,0}
\definecolor{codegray}{rgb}{0.5,0.5,0.5}
\definecolor{codepurple}{rgb}{0.58,0,0.82}
\definecolor{backcolour}{rgb}{0.95,0.95,0.92}

% 设置代码块全局样式
\lstset{
    backgroundcolor=\color{backcolour},   % 背景色
    commentstyle=\color{codegreen},       % 注释颜色
    keywordstyle=\color{blue}\bfseries,   % 关键字颜色与加粗
    numberstyle=\tiny\color{codegray},    % 行号样式
    stringstyle=\color{codepurple},       % 字符串颜色
    basicstyle=\ttfamily\small,           % 基础字体（等宽小号）
    breakatwhitespace=false,         
    breaklines=true,                      % 自动折行（防止代码过长超出边界）
    numbers=left,                         % 行号显示在左侧
    numbersep=5pt,                  
    showspaces=false,                
    showstringspaces=false,
    showtabs=false,                  
    tabsize=4                             % Tab 宽度为 4
}

% 自定义提示框 (用于避坑指南)
\newtcolorbox{tipbox}[1][]{
    colback=blue!5!white,
    colframe=blue!75!black,
    fonttitle=\bfseries,
    title=💡 提示/避坑: #1
}

% 文档信息
\title{\textbf{从源码安装服务器}}
\author{安徽理工大学 RoboCup 实验室}
\date{\today}

% 页眉页脚设置
\usepackage{fancyhdr}
\pagestyle{fancy}
\fancyhf{} % 清空所有默认的页眉页脚设置
% 只设置页脚中央为页码，页眉保持空白
\cfoot{\thepage} 
% 移除页眉那条横线 (Remove the header rule line)
\renewcommand{\headrulewidth}{0pt}

% ============================正文开始==============================================
\begin{document}

\maketitle

本入门指南的源码在\url{https://github.com/LR2933/robocup-3d-getting-started}

欢迎star和pr。
\tableofcontents % 生成目录
\newpage
之前的教程里安装\texttt{rcssservermj}服务器用的都是\texttt{pipx}安装的，这种安装方式简单快速，但是安装的一般都不是最新的版本，所以在这里我们要从源码安装。
\section{拉取源代码}
使用以下命令拉取源代码
\begin{lstlisting}[language=bash]
git clone https://gitlab.com/robocup-sim/rcssservermj.git
\end{lstlisting}

然后当前目录下就会出现一个\texttt{rcssservermj}，使用以下命令切换目录到\texttt{rcssservermj}下
\begin{lstlisting}[language=bash]
cd rcssservermj
\end{lstlisting}

\section{创建虚拟环境}
使用以下命令创建虚拟环境，并且指定python版本是3.12
\begin{lstlisting}[language=bash]
uv venv --python 3.12
\end{lstlisting}

\section{安装依赖}
使用以下命令在当前虚拟环境里安装依赖
\begin{lstlisting}[language=bash]
uv pip install -r requirements.txt
\end{lstlisting}

\section{安装服务器}
使用以下命令在当前环境里安装服务器
\begin{lstlisting}[language=bash]
uv pip install -e .
\end{lstlisting}

\section{启动服务器}
使用以下命令启动服务器
\begin{lstlisting}[language=bash]
uv run rcssservermj
\end{lstlisting}

但是这个命令只能在\texttt{rcssservermj}目录下才可以用，如果不嫌麻烦可以每次都到这个文件里启动

如果嫌麻烦的话可以用下面的方法创建一个别名

\newpage

\section{创建别名}
打开位于家目录下的\texttt{.bashrc}(\url{~/.bashrc})，把下面的命令复制进去(注意\texttt{--directory}后面的是\texttt{rcssservermj}的位置，以你的实际位置为准。如果你是在家目录下\texttt{git clone}的，那么就可以直接复制下面的命令)
\begin{lstlisting}[language=bash]
alias rcss="uv --directory ~/rcssservermj/ run rcssservermj"
\end{lstlisting}

使用以下命令重新加载\texttt{.bashrc}
\begin{lstlisting}[language=bash]
source ~/.bashrc
\end{lstlisting}

然后你就可以在任何地方使用\texttt{rcss}启动服务器了
\bigbreak
如果你觉得敲\texttt{rcss}还是太麻烦了，你甚至可以把\texttt{alias rcss="uv --directory ~/rcssservermj/ run rcssservermj"}里的\texttt{rcss}改成\texttt{r}

这样你就可以在任何地方敲个\texttt{r}启动服务器了
\end{document}

